\section{Graphs covered by four-cycles}\label{sec:squares}
The next bounds impose no restriction on individual block ranks. Their
graph input concerns the edges in the \emph{original} support graph that
belong to a simple four-cycle. Write $G_\square$ for this spanning
subgraph, and $F=E(G)\setminus E(G_\square)$. There is no iterative
edge-deletion process in this definition.

\begin{theorem}[C: a common target for square-covered graphs]\label{thm:square-good}
Let $G$ be a finite simple triangle-free graph of maximum degree three.
If every edge belongs to a four-cycle, then $G$ admits a graph
homomorphism to the circular clique $K_{8/3}$. Consequently every
normalized, monotone, submodular function $h:2^V\to\R$ satisfying
$h(N_G(v)\cup\{v\})=h(N_G(v))$ obeys
\[
 2h(V)\leq\sum_{v\in V}h(\{v\}).
\]
In particular, $G_\square$ satisfies this weighted half bound for every
triangle-free subcubic $G$, whether or not all edges of $G$ lie in squares.
\formal{SquareGood.lean}{Squares.target,Squares.good,squareCore_good}
\end{theorem}
\begin{outline}
Work first in a connected component. Two small configurations force a
closed bipartite component: two vertices with three common neighbors,
and an embedded cube with one vertex deleted. The formal proof treats
both by the same boundary lemma. A bipartite seed has three boundary
vertices, each already adjacent to two internal vertices. A square
through any new edge forces its outside endpoint to touch two boundary
vertices. Any two such pairs intersect, and the degree bound then forces
all outside attachments to be the same vertex. The enlarged component
remains bipartite and maps to an edge of $K_{8/3}$.

Suppose now that no such homomorphism exists. The preceding configurations
are absent. Distinct vertices have at most two common neighbors, and no
vertex can support squares on all three pairs of incident edges: those
three squares would produce a cube cap, or three common neighbors.

An embedded ladder consists of disjoint two-vertex rungs joined by two
rails. A square with an outside attachment extends to a three-rung
ladder; if no attachment exists, the square is a closed bipartite
component. Choose a longest embedded ladder, of at least three rungs.
Each interior vertex already has degree three. At an endpoint, a new
edge must belong to a square through the endpoint's rung. A square
through the last rail instead would give three common neighbors or
three square wedges at the preceding interior vertex. The new square's
other two vertices both lie outside the ladder, by degree saturation.
It would therefore extend the ladder, a contradiction.

Connectedness makes this ladder span the component. Its extra edges
can only form a partial matching between its two end rungs. Complete
that matching: the resulting graph is a closed ladder, with either
parallel or twisted closing rails. Adding two rungs admits a fold onto
the shorter ladder, by inserting a two-edge backtrack. Four finite base
maps suffice: the three-rung twisted ladder and four-rung untwisted
ladder are bipartite; the four-rung twisted ladder is the M\"obius
ladder $M_8$; the five-rung untwisted ladder maps to its five-cycle.
The three-rung untwisted closure is excluded by triangle-freeness.
The Lean kernel checks these finite maps directly.

Label $M_8$ by $\Z/8\Z$, with edge differences $1,-1,4$.
Multiplication by three identifies it with $K_{8/3}$, whose edge
differences are $3,4,5$. Combine the component maps. Finally apply
the fourth paper's formalized homomorphism pullback and its certified
weighted half bound for $K_{8/3}$~\cite{BrendeCurved}. Every edge of
$G_\square$ retains its original witnessing square, so the same theorem
applies to that subgraph.
\end{outline}

For the construction program, squares alone cannot provide a subcubic
support exceeding half capacity, even for arbitrary polymatroid recovery.
The next section accounts for the edges outside those squares without
assuming that deleting them preserves an arbitrary code.
